% page 21 of the facsimile (image p-020 of the scan)
% block 165.1 x 242.0 mm ; left margin 36.9 mm
\pgc{15.240mm}{0.000mm}{
\l{\cel{58}13}
\l{}
\l{}
\l{\sou{akantonar}. (\sou{trans., en)}.\cel{2}(\sou{milit}.) Instalar la armeani en la}
\l{\cel{3}loki diversa di teritorio ube li devas vivar dum kelka}
\l{\cel{3}tempo. - DEFIS.}
\l{}
\l{\sou{akaparar}. (\sou{trans}.)I. (\sou{ulo}) Rezervigar por su, komprar la tota}
\l{\cel{3}quanto de ta o ca varo qua ofresas vende en merkato. - Prenar}
\l{\cel{3}por su ipsa, detrimente di la kunhomi. - II. (\sou{ulu}) (\sou{metaf}.)}
\l{\cel{3}Kaptar (lu) e ne lasar (lu) relatar kun la personi cetera.}
\l{}
\l{}
\l{\sou{akaro}. (\sou{zool.}) Familio de insekti di qui la tipo esas L.\sou{acarus,}}
\l{\cel{3}kun korpo, generale, disko-forma e globatra, sen distingo tre}
\l{\cel{3}preciza di abdomino de cefalotorako. - EFIRSL.}
\l{}
\l{\sou{aklamar}. (\sou{trans}.)Aceptar per klami e krii de entuziasmo - agita}
\l{\cel{3}da plura, multa, personi kune. - DEFIS.}
\l{}
\l{\sou{aklimatar}.\cel{2}(\sou{trans.}) Igar (lu) adaptar su a klimato altra kam}
\l{\cel{3}ta di lua nasko-lando. - (\sou{anke metaf.}) -DEFIRS.}
\l{}
\l{\sou{akneo}.- (\sou{patol}.)Erupto di papuli. - EFIS.}
\l{}
\l{akoluto. (\sou{Eklezio katolika}). Kleriko qua, en kirko, exekutis la}
\l{\cel{3}taski grosiera. - DEFIRS.}
\l{}
\l{\sou{akomodar}.\cel{2}(\sou{trans.}) Aranjar kozo tale ke lu konvenas por ulu.}
\l{\cel{3}- DEFIRS.}
\l{}
\l{\sou{akompanar}.\cel{2}I. Juntar su (ad ulu) por irar adube lu iras. -}
\l{\cel{3}II. (\sou{muziko}) Adjuntar, a la parto precipua quan kantas o pleas}
\l{\cel{3}ulu, partin acesora por igar olta saliar. - EFIRS.}
\l{}
\l{\sou{akonito}. (\sou{bot.)}\cel{1}Planto venenoza, ek la familio \textquotedbl{}renonkulacei\textquotedbl{}.}
\l{\cel{3}- L.\cel{1}\sou{aconitum.}\cel{1}- DEFIRS.}
\l{}
\l{\sou{akonto}.\cel{2}Pago di parto de debo. - DEFIS.}
\l{}
\l{\sou{akoro}. (\sou{bot.}) Genero de aroidei, kun flori hermafrodita, qui}
\l{\cel{3}enkovras la spadico. - L.\cel{1}\sou{acorus calamus}. - L.}
\l{}
\l{\sou{akordar}.\cel{2}(\sou{trans.}) I. (\sou{muziko}) Existigar relato preciza inter}
\l{\cel{3}la soni diversa di muzik-instrumento segun diapazono.}
\l{\cel{3}Regular plura instrumenti muzikala segun la sama diapazono-}
\l{\cel{3}sono. - II. (\sou{gram}.)Igar (per flexioni) vorto relatar kun}
\l{\cel{3}altra vorto di qua lu devas sequar la varii. - DEFIRS.}
\l{}
\l{\sou{akordeono}. Muzik-instrumento, kun klavaro, langeti metala,}
\l{\cel{3}libere movebla, quin vibrigas suflilo. - DEFIRS.}
\l{}
\l{\sou{akra}.\cel{2}I. Di qua la saporo, la odoro, esas forta, quaze acida,}
\l{\cel{3}ed iritas la respir-organi e la guturo. - II. (\sou{metaf})}
\l{\cel{3}(\sou{aludante persono}). Di qua la karaktero esas bitre}
\l{\cel{3}vundiva. - EFIS.}
}
